%-----------------------------------------------------------------------------------------------
\makeatletter
\immediate\write18{datelog > \jobname.info} % site script for $(date '+%Y-%m-%d %Hh%Mm%Ss')
\makeatother
%-----------------------------------------------------------------------------------------------
%-----------------------------------------------------------------------------------------------
\usetheme{Copenhagen}
\usepackage{beamercolorthemeCNThermSci2}
\usefonttheme{serif}
%-----------------------------------------------------------------------------------------------

%-----------------------------------------------------------------------------------------------
%-----------------------------------------------------------------------------------------------
\usetheme{Copenhagen}
\usepackage{beamercolorthemeUTF2}
\usefonttheme{serif}
%-----------------------------------------------------------------------------------------------
\usepackage[utf8]{inputenc}
\usepackage[greek,french,english,brazil]{babel} % last becomes the active one
\usepackage{pslatex}
\usepackage{amssymb,amsmath}
\usepackage{soul}
\usepackage[squaren,Gray,cdot]{SIunits}
\usepackage[nice]{nicefrac}
\usepackage{tikz}
\usepackage{amscd}
\usepackage{stmaryrd}
\usepackage{scalerel}
\usepackage{xspace}
%-----------------------------------------------------------------------------------------------


%-----------------------------------------------------------------------------------------------
%-----------------------------------------------------------------------------------------------
% Mathematical
%-----------------------------------------------------------------------------------------------
\newcommand{\vet}[1]{\underline{{#1}}}
\newcommand{\mat}[1]{\underline{\underline{{#1}}}}
\newcommand{\cub}[1]{\underline{\underline{\underline{{#1}}}}}
\newcommand{\eqdef}{{\ensuremath\stackrel{\text{\tiny def}}{=}}}
%-----------------------------------------------------------------------------------------------
% Linguistic
%-----------------------------------------------------------------------------------------------
\newcommand{\GRtxt}[1]{\begin{otherlanguage}{greek}{{#1}}\end{otherlanguage}}
\newcommand{\FRtxt}[1]{\begin{otherlanguage}{french}{{#1}}\end{otherlanguage}}
%-----------------------------------------------------------------------------------------------
% Presentation
%-----------------------------------------------------------------------------------------------
\newcommand{\BkgImgH}[1]{% Places an image centered on the slide background filling the height
    \usebackgroundtemplate{\parbox{\paperwidth}{%
        \vspace*{1sp}\centering\includegraphics[height=\paperheight]{{#1}}
}}}
\newcommand{\BkgImgW}[1]{% Places an image centered on the slide background filling the width
    \usebackgroundtemplate{\parbox{\paperwidth}{%
        \vspace*{1sp}\centering\includegraphics[width=\paperwidth]{{#1}}
}}}
\newcommand{\ArtEndH}[3]{% Transitions to plain image (last) slide: #1:prefix #2,#3:extensions
    \BkgImgH{root/../art/#1.#2}
    \frame<handout:0>[plain]{%
        \transdissolve\vspace*{72mm}\color{white}\scriptsize\bf\input{root/../art/#1.#3}}
    \usebackgroundtemplate{\mbox{~}}
}
\newcommand{\ArtEndW}[3]{% Transitions to plain image (last) slide: #1:prefix #2,#3:extensions
    \BkgImgW{root/../art/#1.#2}
    \frame<handout:0>[plain]{%
        \transdissolve\vspace*{72mm}\color{white}\scriptsize\bf\input{root/../art/#1.#3}}
    \usebackgroundtemplate{\mbox{~}}
}
\newcommand{\ImgColW}[3]{% Inserts a full-width image in a column
    \includegraphics[width=\columnwidth]{root/../art/#1.#2}\\[-0.5\baselineskip]
    \parbox{\columnwidth}{\tiny\hfill\scalebox{0.85}{\input{root/../art/#1.#3}}}
}
\newcommand{\txtpic}[1]{%
    \fcolorbox{lightgray}{white!90!black}{{#1}} 
}
%-----------------------------------------------------------------------------------------------


%-----------------------------------------------------------------------------------------------
\title{A.07.01 -- Relações de Propriedades Termodinâmicas}
\subtitle{Relações Diferenciais Parciais e de Maxwell}
\author{Prof.~C.~Naaktgeboren, PhD}
\date{{\scriptsize\tt%
    \includegraphics[height=6.0mm]{cc/by-nc-nd-88x31.pdf}\\[\smallskipamount]
    https://github.com/CNThermSci/ApplThermSci\\
    Compiled on \input{\jobname.info}
}}
%-----------------------------------------------------------------------------------------------
\begin{document}
%-----------------------------------------------------------------------------------------------
\logo{%
    \parbox{158mm}{% There's a 1mm gap on each side of the 160mm x 90mm slide logo line
        \mode<beamer>{
            \includegraphics[height=6.0mm]{root/00-res/UTFPR/UTFPR-logo-D.pdf}\hfill%
            \includegraphics[height=9.0mm]{root/00-res/logo/CNThermSci-logo-A.pdf}%
        }
        \mode<handout>{
            \includegraphics[height=6.0mm]{root/00-res/UTFPR/UTFPR-logo-W.pdf}\hfill%
            \includegraphics[height=9.0mm]{root/00-res/logo/CNThermSci-logo-W.pdf}%
        }
    }
} % The (delineated, alpha), or washed-out logos
%-----------------------------------------------------------------------------------------------
\frame{\titlepage}
%-----------------------------------------------------------------------------------------------

%-----------------------------------------------------------------------------------------------
\frame{\tableofcontents}
%-----------------------------------------------------------------------------------------------

%-----------------------------------------------------------------------------------------------
\section{Relações Diferenciais Parciais}
%-----------------------------------------------------------------------------------------------

%-----------------------------------------------------------------------------------------------
\subsection{Função de Duas Variáveis Independentes}
%-----------------------------------------------------------------------------------------------

    % !j 96 -i8
    %-------------------------------------------------------------------------------------------
    \begin{frame}{Diferencial Total}\vspace*{-2em}
        Seja \alert{$z\!:\!z(x, y)$}, com \alert{$x$} e \alert{$y$} contínuos e diferenciáveis,
        então:

        \begin{equation*}
            \alert{dz =
                \left(
                    \frac{\partial z}{\partial x}
                \right)\,dx +
                \left(
                    \frac{\partial z}{\partial y}
                \right)\,dy},
        \end{equation*}\vspace*{\medskipamount}

        As derivadas parciais são efetuadas mantendo as \alert{demais variáveis constantes}:

        A notação abaixo --- como em \alert{$c_p$} e \alert{$c_v$}, por exemplo ---
        \alert{explicita} esta condição:

        \begin{equation*}
            \alert{dz =
                \left(
                    \frac{\partial z}{\partial x}
                \right)_y\,dx +
                \left(
                    \frac{\partial z}{\partial y}
                \right)_x\,dy}.
        \end{equation*}
    \end{frame}
    %-------------------------------------------------------------------------------------------

    % !j 96 -i8
    %-------------------------------------------------------------------------------------------
    \begin{frame}{Relação de Derivadas Parciais Cruzadas}\vspace*{-2em}
        Escrevendo-se \alert{$dz = Mdx + Ndy$}, com
        \begin{equation*}
            M =
                \left(
                    \frac{\partial z}{\partial x}
                \right)_y
            \qquad\mbox{e}\qquad
            N =
                \left(
                    \frac{\partial z}{\partial y}
                \right)_x.
        \end{equation*}

        Então, fazendo:
        \begin{align*}
            \left(
                \frac{\partial^2 z}{\partial x\partial y}
            \right) &=
            \left(
                \frac{\partial^2 z}{\partial y\partial x}
            \right),
            \qquad\mbox{resulta em:}
            \\[\bigskipamount]
            \alert{\left(
                \frac{\partial M}{\partial y}
            \right)_x} &\alert{=
            \left(
                \frac{\partial N}{\partial x}
            \right)_y}.
        \end{align*}
    \end{frame}
    %-------------------------------------------------------------------------------------------

%-----------------------------------------------------------------------------------------------
\section{Tópicos de Leitura}
%-----------------------------------------------------------------------------------------------

    %------------------------------------------------------------------------------------------
    \begin{frame}[allowframebreaks]{Tópicos de Leitura}
        \begin{thebibliography}{Çengel, Y.~A., 2013}
            \bibitem[Çengel, Y.~A., 2013]{2013-CengelYA+BolesMA-AMGH}
                Çengel, Y.~A. e Boles, M.~A.
                \newblock{%
                    {\em Termodinâmica $7^\mathrm{a}\!$ Edição\/}.
                    \alert{Seções~0-0 a 0-0.}
                }
                \newblock{\footnotesize AMGH. Porto Alegre. ISBN 978-85-8055-200-3.}
        \end{thebibliography}
    \end{frame}
    %------------------------------------------------------------------------------------------

    % Finishes with stunning image, with credit
    \ArtEndWSE{pexels-emiliano-arano-1295138}{jpg}{txt}

%-----------------------------------------------------------------------------------------------
\end{document}
%-----------------------------------------------------------------------------------------------

